% 三层高斯金字塔与拉普拉斯金字塔的构造关系示意；方块边长按每层减半绘制。宽约 104mm。
\begin{tikzpicture}[x=1mm,y=1mm,font=\figurefont\small,
  lvl/.style={draw=FigureBlue,fill=FigureBlueFill,line width=.6pt},
  lap/.style={draw=FigureYellow,fill=FigureYellowFill,line width=.6pt},
  link/.style={-{Stealth[length=2mm]},draw=FigureStroke,line width=.7pt},
  lab/.style={font=\figurefont\footnotesize,text=FigureMuted}]
  \draw[lvl] (0,0) rectangle (24,24);   \node at (12,12) {$G_0$};
  \draw[lvl] (32,6) rectangle (44,18);  \node at (38,12) {$G_1$};
  \draw[lvl] (64,9) rectangle (70,15);  \node at (67,12) {$G_2$};
  \draw[link] (24.5,12)--(31.5,12);
  \draw[link] (44.5,12)--(63.5,12);
  \node[lab] at (28,27) {模糊、隔点取样};
  \draw[lap] (0,-34) rectangle (24,-10); \node at (12,-22) {$L_0$};
  \draw[lap] (32,-28) rectangle (44,-16); \node at (38,-22) {$L_1$};
  \draw[lvl] (64,-25) rectangle (70,-19); \node at (67,-22) {$G_2$};
  \draw[link] (12,-.5)--(12,-9.5) node[midway,right,lab] {减去扩展的$G_1$};
  \draw[link] (38,5.5)--(38,-15.5) node[midway,right,lab] {减去扩展的$G_2$};
  \draw[link] (67,8.5)--(67,-18.5) node[midway,right,lab] {保留};
  \node[anchor=west] at (80,12) {高斯金字塔};
  \node[anchor=west] at (80,-22) {拉普拉斯金字塔};
\end{tikzpicture}
